Before stating the definition of pre-analytic rings, we first make the following construction. Let $\mathsf{FF}$ denote the full subcategory of $\mathsf{Fun}(\Delta^1,\widehat{\mathsf{Cat}})$ consisting of fully faithful functors, and consider the pullback diagram
% https://q.uiver.app/#q=WzAsNCxbMCwwLCJcXG1hdGhzZntQfV97XFxsYW1iZGF9Il0sWzAsMSwiXFxtYXRoc2Z7Q29uZH1fe1xcbGFtYmRhfShcXG1hdGhzZntSaW5nfSkiXSxbMiwxLCJcXHdpZGVoYXR7XFxtYXRoc2Z7Q2F0fX1eXFxtYXRocm17b3B9Il0sWzIsMCwiXFxtYXRoc2Z7RkZ9XlxcbWF0aHJte29wfSJdLFszLDIsInQiXSxbMSwyLCJcXG1hdGhzZntEfSgtKSJdLFswLDFdLFswLDNdXQ==
\[\begin{tikzcd}[ampersand replacement=\&]
	{\mathsf{P}_{\lambda}} \&\& {\mathsf{FF}^\mathrm{op}} \\
	{\mathsf{Cond}_{\lambda}(\mathsf{Ring})} \&\& {\widehat{\mathsf{Cat}}^\mathrm{op},}
	\arrow[from=1-1, to=1-3]
	\arrow[from=1-1, to=2-1]
	\arrow["t", from=1-3, to=2-3]
	\arrow["{\mathsf{D}(-)}", from=2-1, to=2-3]
\end{tikzcd}\]
where $t \colon \mathsf{FF} \to \widehat{\mathsf{Cat}}$ sends a functor to its target.

Since $t$ is an isofibration, \cite[Tag \href{https://kerodon.net/tag/033T}{033T}]{kerodon} provides an explicit description of $\mathsf{P}_\lambda$. Its objects are pairs $A=(A^\triangleright, \mathsf{D}(A))$, where $A^\triangleright$ is a $\lambda$-light condensed ring and $\mathsf{D}(A)$ is a full subcategory of $\mathsf{D}(A^\triangleright)$. A morphism $A \to B$ in $\mathsf{P}_\lambda$ consists of a morphism of $\lambda$-light condensed rings $\varphi\colon A^\triangleright \to B^\triangleright$ such that the associated restriction-of-scalars functor
\[\varphi_*\colon \mathsf{D}(B^\triangleright) \to \mathsf{D}(A^\triangleright)\]
carries $\mathsf{D}(B)$ into $\mathsf{D}(A)$. 

\begin{definition}\label{def panring}
A \emph{$\lambda$-light pre-analytic ring} $A$ is a pair $(A^\triangleright,\mathsf{D}(A))$ consisting of a $\lambda$-light condensed ring $A^\triangleright$ and a full subcategory $\mathsf{D}(A)$ of $\mathsf{D}(A^\triangleright)$ such that:
\begin{enumerate}[label=(\roman*)]
    \item The full subcategory $\mathsf{D}(A)\subset \mathsf{D}(A^\triangleright)$ is closed under small limits and colimits (hence the inclusion $\iota_A\colon \mathsf{D}(A)\hookrightarrow\mathsf{D}(A^\triangleright)$ admits a left adjoint, denoted $-\otimes_{A^\triangleright}A$);
    \item For every $M \in \mathsf{D}(A^\triangleright)$ and $N \in \mathsf{D}(A)$, the object $\underline{\operatorname{Hom}}(M,N)$ lies in $\mathsf{D}(A)$;
    \item The composition $L_A\colonequals\iota_A \circ \bigl(-\otimes_{A^\triangleright}A\bigr)$ preserves connective objects.
\end{enumerate}
We define $\mathsf{PAnRing}_{\lambda}$ to be the full subcategory of $\mathsf{P}_{\lambda}$ spanned by the $\lambda$-light pre-analytic rings.
\end{definition}

\begin{theorem}\label{PAnRing is presentable}
    The category $\mathsf{PAnRing}_\lambda$ is $\mu$-presentable, for $\mu=((2^{<\lambda})^{\aleph_0})^+$.
\end{theorem}

The proof of this theorem requires us to reinterpret each of the three conditions in Definition~\ref{def panring}. We begin with the third, which leads us to consider a variant $\mathsf{P}_\lambda^\mathrm{cn}$ of $\mathsf{P}_\lambda$, obtained by replacing the functor $\mathsf{D}(-)$ in its definition with
\[\mathsf{D}(-)_{\geq 0}\colon \mathsf{Cond}_\lambda(\mathsf{Ring})\to \widehat{\mathsf{Cat}}^\mathrm{op}.\]
In the same spirit, we introduce the following connective analogue of Definition~\ref{def panring}.

\begin{definition}\label{def connective panring}
We denote by $\mathsf{PAnRing}^\mathrm{cn}_\lambda$ the full category of $\mathsf{P}_\lambda^\mathrm{cn}$ spanned by pairs $A=(A^{\triangleright},\mathsf{D}(A)_{\geq0})$, consisting of a $\lambda$-light condensed ring $A^{\triangleright}$ and a full subcategory $\mathsf{D}(A)_{\geq0}$ of $\mathsf{D}(A^{\triangleright})_{\geq0}$ such that: 
    \begin{enumerate}[label=(\roman*)]
        \item The full subcategory $\mathsf{D}(A)_{\geq0}\subset\mathsf{D}(A^{\triangleright})_{\geq0}$ is closed under small limits and colimits (hence the inclusion admits a left adjoint, which we still denote by $-\otimes_{A^{\triangleright}}A$);
        \item For every $M \in \mathsf{D}(A^\triangleright)_{\geq0}$ and $N \in \mathsf{D}(A)_{\geq0}$, the object $\tau_{\geq 0}\underline{\operatorname{Hom}}(M,N)$ lies in $\mathsf{D}(A)_{\geq0}$.
    \end{enumerate}
\end{definition}

There is a natural comparison between $\mathsf{PAnRing}_\lambda$ and $\mathsf{PAnRing}_\lambda^{\mathrm{cn}}$. To construct it, let $\mathsf{FF}^{\mathsf{R}}$ denote the full subcategory of $\mathsf{Fun}(\Delta^1,\mathsf{Pr}^{\mathsf{R}})$ spanned by the fully faithful functors. The reflection theorem implies that the composite $\mathsf{PAnRing}^\mathrm{cn}_\lambda\hookrightarrow \mathsf{P}^\mathrm{cn}_\lambda\to \mathsf{FF}^\mathrm{op}$ factors as
% https://q.uiver.app/#q=WzAsNCxbMCwwLCJcXG1hdGhzZntQQW5SaW5nfV9cXGxhbWJkYSJdLFswLDEsIlxcbWF0aHNme1B9X1xcbGFtYmRhIl0sWzIsMCwiXFxtYXRoc2Z7RkZ9XntcXG1hdGhzZntSfSxcXG1hdGhybXtvcH19Il0sWzIsMSwiXFxtYXRoc2Z7RkZ9XlxcbWF0aHJte29wfSwiXSxbMCwxLCIiLDAseyJzdHlsZSI6eyJ0YWlsIjp7Im5hbWUiOiJob29rIiwic2lkZSI6InRvcCJ9fX1dLFsxLDNdLFsyLDMsIiIsMix7InN0eWxlIjp7InRhaWwiOnsibmFtZSI6Imhvb2siLCJzaWRlIjoidG9wIn19fV0sWzAsMiwiIiwyLHsic3R5bGUiOnsiYm9keSI6eyJuYW1lIjoiZGFzaGVkIn19fV1d
\[\begin{tikzcd}[ampersand replacement=\&]
	{\mathsf{PAnRing}^\mathrm{cn}_\lambda} \&\& {\mathsf{FF}^{\mathsf{R},\mathrm{op}}} \\
	{\mathsf{P}^\mathrm{cn}_\lambda} \&\& {\mathsf{FF}^\mathrm{op}.}
	\arrow[dashed, from=1-1, to=1-3]
	\arrow[hook, from=1-1, to=2-1]
	\arrow[hook, from=1-3, to=2-3]
	\arrow[from=2-1, to=2-3]
\end{tikzcd}\]

Now let $A^\triangleright$ be a $\lambda$-light condensed ring. By Lemma~\ref{t-structures on derived categories}, there is an equivalence of categories $\mathsf{D}(A^\triangleright)\simeq \mathsf{D}(A^\triangleright)_{\geq 0}\otimes \mathsf{Sp}$. It follows from the universal property of the fiber product that these equivalences assemble into the dashed functor $\Psi_\lambda$ in the diagram below:
% https://q.uiver.app/#q=WzAsOSxbMSwxLCJcXG1hdGhzZntQfV9cXGxhbWJkYSJdLFszLDEsIlxcbWF0aHNme0ZGfV5cXG1hdGhybXtvcH0sIl0sWzEsMiwiXFxtYXRoc2Z7Q29uZH1fXFxsYW1iZGEoXFxtYXRoc2Z7UmluZ30pIl0sWzMsMiwiXFx3aWRlaGF0e1xcbWF0aHNme0NhdH19XlxcbWF0aHJte29wfSJdLFswLDEsIlxcbWF0aHNme1B9X1xcbGFtYmRhXlxcbWF0aHJte2NufSJdLFswLDIsIlxcbWF0aHNme0NvbmR9X1xcbGFtYmRhKFxcbWF0aHNme1Jpbmd9KSJdLFswLDAsIlxcbWF0aHNme1BBblJpbmd9X1xcbGFtYmRhXlxcbWF0aHJte2NufSJdLFsxLDAsIlxcbWF0aHNme0ZGfV57XFxtYXRoc2Z7Un0sXFxtYXRocm17b3B9fSJdLFszLDAsIlxcbWF0aHNme0ZGfV57XFxtYXRoc2Z7Un0sXFxtYXRocm17b3B9fSJdLFswLDFdLFswLDJdLFsyLDMsIlxcbWF0aHNme0R9KC0pIl0sWzEsMywidCJdLFs1LDIsIiIsMCx7InN0eWxlIjp7ImhlYWQiOnsibmFtZSI6Im5vbmUifX19XSxbNCw1XSxbNiw0LCIiLDAseyJzdHlsZSI6eyJ0YWlsIjp7Im5hbWUiOiJob29rIiwic2lkZSI6InRvcCJ9fX1dLFs2LDddLFs3LDgsIi1cXG90aW1lc1xcbWF0aHNme1NwfSJdLFs4LDEsIiIsMix7InN0eWxlIjp7InRhaWwiOnsibmFtZSI6Imhvb2siLCJzaWRlIjoidG9wIn19fV0sWzYsMCwiIiwxLHsic3R5bGUiOnsiYm9keSI6eyJuYW1lIjoiZGFzaGVkIn19fV1d
\[\begin{tikzcd}[ampersand replacement=\&]
	{\mathsf{PAnRing}_\lambda^\mathrm{cn}} \& {\mathsf{FF}^{\mathsf{R},\mathrm{op}}} \&\& {\mathsf{FF}^{\mathsf{R},\mathrm{op}}} \\
	{\mathsf{P}_\lambda^\mathrm{cn}} \& {\mathsf{P}_\lambda} \&\& {\mathsf{FF}^\mathrm{op},} \\
	{\mathsf{Cond}_\lambda(\mathsf{Ring})} \& {\mathsf{Cond}_\lambda(\mathsf{Ring})} \&\& {\widehat{\mathsf{Cat}}^\mathrm{op}.}
	\arrow[from=1-1, to=1-2]
	\arrow[hook, from=1-1, to=2-1]
	\arrow["{\Psi_\lambda}"', dashed, from=1-1, to=2-2]
	\arrow["{-\otimes\mathsf{Sp}}", from=1-2, to=1-4]
	\arrow[hook, from=1-4, to=2-4]
	\arrow[from=2-1, to=3-1]
	\arrow[from=2-2, to=2-4]
	\arrow[from=2-2, to=3-2]
	\arrow["t", from=2-4, to=3-4]
	\arrow[equals, from=3-1, to=3-2]
	\arrow["{\mathsf{D}(-)}", from=3-2, to=3-4]
\end{tikzcd}\]

\begin{lemma}\label{preanalytic rings are determined by connective halves}
The functor $\Psi_\lambda\colon \mathsf{PAnRing}_\lambda^{\mathrm{cn}}\to \mathsf{P}_\lambda$ is fully faithful, with essential image $\mathsf{PAnRing}_\lambda$. In particular, it induces an equivalence $\mathsf{PAnRing}_\lambda^{\mathrm{cn}}\simeq \mathsf{PAnRing}_\lambda$.
\end{lemma}

\begin{proof}
We first outline the argument, before giving the details of each step. After observing that the essential image of $\Psi_\lambda$ is contained in $\mathsf{PAnRing}_\lambda$, we obtain a commutative diagram
% https://q.uiver.app/#q=WzAsMyxbMiwwLCJcXG1hdGhzZntQQW5SaW5nfV9cXGxhbWJkYSJdLFsxLDEsIlxcbWF0aHNme0NvbmR9X1xcbGFtYmRhKFxcbWF0aHNme1Jpbmd9KSwiXSxbMCwwLCJcXG1hdGhzZntQQW5SaW5nfV9cXGxhbWJkYV5cXG1hdGhybXtjbn0iXSxbMCwxLCJVIl0sWzIsMCwiXFxQc2lfXFxsYW1iZGEiXSxbMiwxLCJVXlxcbWF0aHJte2NufSIsMl1d
\[\begin{tikzcd}[ampersand replacement=\&,column sep = tiny]
	{\mathsf{PAnRing}_\lambda^\mathrm{cn}} \&\& {\mathsf{PAnRing}_\lambda} \\
	\& {\mathsf{Cond}_\lambda(\mathsf{Ring}),}
	\arrow["{\Psi_\lambda}", from=1-1, to=1-3]
	\arrow["{U^\mathrm{cn}}"', from=1-1, to=2-2]
	\arrow["U", from=1-3, to=2-2]
\end{tikzcd}\]
where $U^\mathrm{cn}$ and $U$ forget the categories of complete modules. We then show that both $U^\mathrm{cn}$ and $U$ are cocartesian fibrations and that $\Psi_\lambda$ preserves cocartesian maps. According to \cite[Tag~\href{https://kerodon.net/tag/028B}{028B}]{kerodon}, it therefore suffices to show that $\Psi_\lambda$ induces an equivalence on each fiber. On fibers, this reduces to the classical fact that pre-analytic ring structures can be encoded in terms of abelian categories.

\textsc{Step I: The essential image of $\Psi_\lambda$ is contained in $\mathsf{PAnRing}_\lambda$.} Let $(A^\triangleright,\mathsf{D}(A)_{\geq 0})$ be an object of $\mathsf{PAnRing}_\lambda^\mathrm{cn}$. Since stabilization is a 2-functor \cite[Rem.~6.1.7]{CLL}, we obtain an adjunction
% https://q.uiver.app/#q=WzAsMixbMCwwLCJcXG1hdGhzZntEfShBKV97XFxnZXEgMH1cXG90aW1lcyBcXG1hdGhzZntTcH0iXSxbMiwwLCJcXG1hdGhzZntEfShBXlxcdHJpYW5nbGVyaWdodClfe1xcZ2VxIDB9XFxvdGltZXMgXFxtYXRoc2Z7U3B9LiJdLFswLDEsIiIsMCx7ImN1cnZlIjotM31dLFsxLDAsIiIsMCx7ImN1cnZlIjotMywic3R5bGUiOnsidGFpbCI6eyJuYW1lIjoiaG9vayIsInNpZGUiOiJib3R0b20ifX19XSxbMiwzLCIiLDIseyJsZXZlbCI6MSwic3R5bGUiOnsibmFtZSI6ImFkanVuY3Rpb24ifX1dXQ==
\[\begin{tikzcd}[ampersand replacement=\&]
	{\mathsf{D}(A^\triangleright)_{\geq 0}\otimes \mathsf{Sp}} \&\& {\mathsf{D}(A)_{\geq 0}\otimes \mathsf{Sp}.}
	\arrow[""{name=0, anchor=center, inner sep=0}, "{(- \otimes_{A^\triangleright}A)\otimes \mathsf{Sp}}", bend left, from=1-1, to=1-3]
	\arrow[""{name=1, anchor=center, inner sep=0}, bend left, hook', from=1-3, to=1-1]
	\arrow["\dashv"{anchor=center, rotate=-90}, draw=none, from=0, to=1]
\end{tikzcd}\]
In particular, the full subcategory $\mathsf{D}(A)_{\geq 0}\otimes \mathsf{Sp}\subseteq \mathsf{D}(A^\triangleright)_{\geq 0}\otimes \mathsf{Sp}\simeq \mathsf{D}(A^\triangleright)$ is stable under small limits. On the other hand, the inclusion $\mathsf{D}(A)_{\geq 0}\hookrightarrow \mathsf{D}(A^\triangleright)_{\geq 0}$ preserves small colimits and therefore admits a right adjoint. Stabilizing this adjunction shows that $\mathsf{D}(A)_{\geq 0}\otimes \mathsf{Sp}\subseteq \mathsf{D}(A^\triangleright)$ is also closed under small colimits.

By Proposition~\ref{enhanced reflection theorem}, the completion functor
\[-\otimes_{A^\triangleright}A\colon \mathsf{D}(A^\triangleright)_{\geq 0}\to \mathsf{D}(A)_{\geq 0}\]
endows $\mathsf{D}(A)_{\geq 0}$ with a symmetric monoidal structure for which this functor is symmetric monoidal. Moreover, Proposition~\ref{enhanced reflection theorem} shows that it defines a morphism in $\mathsf{CAlg}(\mathsf{Pr}^\mathsf{L}_{\aleph_1})$. Applying the argument of Proposition~\ref{D is a map in Pr_k}, we obtain the corresponding statement for
\[(-\otimes_{A^\triangleright}A)\otimes\mathsf{Sp}.\]
In particular, this functor is \emph{compatible with the symmetric monoidal structure} in the sense of \cite[Def.~2.2.1.6]{higheralgebra}. Condition~(ii) of Definition~\ref{def panring} now follows from \cite[Lem.~7.5]{reflectionprinciple}.

Finally, condition~(iii) of Definition~\ref{def panring} follows by applying \cite[Prop.~C.3.1.1]{lurie2018spectral} to the inclusion $\mathsf{D}(A)_{\geq 0}\hookrightarrow\mathsf{D}(A^\triangleright)_{\geq 0}$ and to the completion functor $-\otimes_{A^\triangleright}A$. 

\textsc{Step II: $U^\mathrm{cn}$ and $U$ are cocartesian fibrations.} We treat $U$, as the same argument applies to $U^\mathrm{cn}$. Let $A=(A^\triangleright,\mathsf{D}(A))$ be a $\lambda$-light pre-analytic ring, and let $\varphi\colon A^\triangleright\to B^\triangleright$ be a morphism of $\lambda$-light condensed rings. Write $B^\triangleright_{/A}$ for the induced pre-analytic ring structure on $B^\triangleright$ \cite[Lem.~4.1.9]{camargo2026notessolidgeometry}. Thus, we have a pushout square of $\lambda$-light pre-analytic rings
    % https://q.uiver.app/#q=WzAsNCxbMCwwLCJBXlxcdHJpYW5nbGVyaWdodCJdLFswLDEsIkJeXFx0cmlhbmdsZXJpZ2h0Il0sWzEsMCwiQSJdLFsxLDEsIkJeXFx0cmlhbmdsZXJpZ2h0L0EuIl0sWzAsMl0sWzIsMywiZiJdLFswLDEsImZeXFx0cmlhbmdsZXJpZ2h0IiwyXSxbMSwzXV0=
\[\begin{tikzcd}[ampersand replacement=\&]
	{A^\triangleright} \& A \\
	{B^\triangleright} \& {B^\triangleright_{/A}.}
	\arrow[from=1-1, to=1-2]
	\arrow["{\varphi}"', from=1-1, to=2-1]
	\arrow["\overline{\varphi}", from=1-2, to=2-2]
	\arrow[from=2-1, to=2-2]
\end{tikzcd}\]
The functor $U$ maps $\overline{\varphi}\colon A\to B^\triangleright_{/A}$ to $\varphi$, so it suffices to prove that $\overline{\varphi}$ is a $U$-cocartesian morphism.

For any $\lambda$-light pre-analytic ring $C$, the universal property of the pushout gives a cartesian square
% https://q.uiver.app/#q=WzAsNCxbMSwxLCJcXG9wZXJhdG9ybmFtZXtNYXB9X3tcXG1hdGhzZntQQW5SaW5nfX0oQV5cXHRyaWFuZ2xlcmlnaHQsQykuIl0sWzAsMSwiXFxvcGVyYXRvcm5hbWV7TWFwfV97XFxtYXRoc2Z7UEFuUmluZ319KEJeXFx0cmlhbmdsZXJpZ2h0LEMpIl0sWzEsMCwiXFxvcGVyYXRvcm5hbWV7TWFwfV97XFxtYXRoc2Z7UEFuUmluZ319KEEsQykiXSxbMCwwLCJcXG9wZXJhdG9ybmFtZXtNYXB9X3tcXG1hdGhzZntQQW5SaW5nfX0oQl5cXHRyaWFuZ2xlcmlnaHQvQSxDKSJdLFsyLDBdLFsxLDBdLFszLDFdLFszLDJdXQ==
\[\begin{tikzcd}[ampersand replacement=\&]
	{\operatorname{Map}_{\mathsf{PAnRing}_\lambda}(B^\triangleright_{/A},C)} \& {\operatorname{Map}_{\mathsf{PAnRing}_\lambda}(A,C)} \\
	{\operatorname{Map}_{\mathsf{PAnRing}_\lambda}(B^\triangleright,C)} \& {\operatorname{Map}_{\mathsf{PAnRing}_\lambda}(A^\triangleright,C).}
	\arrow[from=1-1, to=1-2]
	\arrow[from=1-1, to=2-1]
	\arrow[from=1-2, to=2-2]
	\arrow[from=2-1, to=2-2]
\end{tikzcd}\]
By \cite[Lem.~4.1.5]{camargo2026notessolidgeometry}, we have
\[\operatorname{Map}_{\mathsf{PAnRing}_\lambda}(A^\triangleright,C)\simeq \operatorname{Map}_{\mathsf{Cond}_\lambda(\mathsf{Ring})}(A^\triangleright,C^\triangleright),\]
and similarly for $B^\triangleright$. Hence, the square above identifies with the cartesian square
% https://q.uiver.app/#q=WzAsNCxbMSwxLCJcXG9wZXJhdG9ybmFtZXtNYXB9X3tcXG1hdGhzZntDb25kfV9cXGxhbWJkYShcXG1hdGhzZntSaW5nfSl9KEFeXFx0cmlhbmdsZXJpZ2h0LEMpLCJdLFswLDEsIlxcb3BlcmF0b3JuYW1le01hcH1fe1xcbWF0aHNme0NvbmR9X1xcbGFtYmRhKFxcbWF0aHNme1Jpbmd9KX0oQl5cXHRyaWFuZ2xlcmlnaHQsQykiXSxbMSwwLCJcXG9wZXJhdG9ybmFtZXtNYXB9X3tcXG1hdGhzZntQQW5SaW5nfX0oQSxDKSJdLFswLDAsIlxcb3BlcmF0b3JuYW1le01hcH1fe1xcbWF0aHNme1BBblJpbmd9fShCXlxcdHJpYW5nbGVyaWdodC9BLEMpIl0sWzIsMF0sWzEsMF0sWzMsMV0sWzMsMl1d
\[\begin{tikzcd}[ampersand replacement=\&]
	{\operatorname{Map}_{\mathsf{PAnRing}_\lambda}(B^\triangleright_{/A},C)} \& {\operatorname{Map}_{\mathsf{PAnRing}_\lambda}(A,C)} \\
	{\operatorname{Map}_{\mathsf{Cond}_\lambda(\mathsf{Ring})}(B^\triangleright,C)} \& {\operatorname{Map}_{\mathsf{Cond}_\lambda(\mathsf{Ring})}(A^\triangleright,C),}
	\arrow[from=1-1, to=1-2]
	\arrow[from=1-1, to=2-1]
	\arrow[from=1-2, to=2-2]
	\arrow[from=2-1, to=2-2]
\end{tikzcd}\]
which proves that $\overline{\varphi}$ is $U$-cocartesian and completes the proof.

\textsc{Step III: The functor $\Psi_\lambda$ preserves cocartesian morphisms.} Running the preceding argument in reverse, we see that a morphism $A\to B$ in $\mathsf{PAnRing}_\lambda$ (resp.~in $\mathsf{PAnRing}^\mathrm{cn}_\lambda$) is $U$-cocartesian (resp.~$U^\mathrm{cn}$-cocartesian) if and only if the diagram
% https://q.uiver.app/#q=WzAsNCxbMCwwLCJBXlxcdHJpYW5nbGVyaWdodCJdLFsxLDAsIkJeXFx0cmlhbmdsZXJpZ2h0Il0sWzAsMSwiQSJdLFsxLDEsIkIiXSxbMCwyXSxbMCwxXSxbMiwzXSxbMSwzXV0=
\[\begin{tikzcd}[ampersand replacement=\&]
	{A^\triangleright} \& {B^\triangleright} \\
	A \& B
	\arrow[from=1-1, to=1-2]
	\arrow[from=1-1, to=2-1]
	\arrow[from=1-2, to=2-2]
	\arrow[from=2-1, to=2-2]
\end{tikzcd}\]
is a pushout, or equivalently, if and only if $B$ carries the pre-analytic ring structure on $B^\triangleright$ induced from $A$.

It therefore remains to show that $\Psi_\lambda$ preserves induced pre-analytic ring structures. Let $A=(A^\triangleright,\mathsf{D}(A)_{\geq 0})$ be an object of $\mathsf{PAnRing}^\mathrm{cn}_\lambda$ and let $\varphi\colon A^\triangleright\to B^\triangleright$ be a morphism of $\lambda$-light condensed rings. By \cite[Lem.~4.1.9]{camargo2026notessolidgeometry}, the pushout $B^\triangleright_{A/}$ defining the induced pre-analytic ring structure exists, and its connective derived category admits the explicit description
\[\mathsf{D}(B^\triangleright_{A/})_{\geq 0}\simeq \{ M\in \mathsf{D}(B^\triangleright)_{\geq 0} \mid \varphi_*M\in \mathsf{D}(A)_{\geq 0}\}.\]
Thus, it suffices to verify that the stabilization $\mathsf{D}(B^\triangleright_{A/})_{\geq 0}\otimes\mathsf{Sp}$ admits the analogous description in terms of $\mathsf{D}(B^\triangleright)\simeq \mathsf{D}(B^\triangleright)_{\geq0}\otimes \mathsf{Sp}$ and $\mathsf{D}(A)_{\geq0}\otimes \mathsf{Sp}$.

By \cite[Prop.~4.1.8]{camargo2026notessolidgeometry}, an object $M\in\mathsf{D}(B^\triangleright)\simeq\mathsf{D}(B^\triangleright)_{\geq 0}\otimes\mathsf{Sp}$ belongs to $\mathsf{D}(B^\triangleright_{A/})_{\geq 0}\otimes\mathsf{Sp}$ if and only if $\pi_i(M)$ does for every $i\in\mathbb{Z}$. Since $\pi_i(M)$ is static, this is equivalent to requiring that $\pi_i(M)$ belong to $\mathsf{D}(B^\triangleright_{A/})_{\geq 0}$ for every $i\in\mathbb{Z}$. By the description above, the latter condition amounts to
\[\varphi_*\pi_i(M)\in\mathsf{D}(A)_{\geq 0}\qquad\text{for all }i\in\mathbb{Z}.\]
Since restriction of scalars is $t$-exact, we have $\varphi_*\pi_i(M)\simeq\pi_i(\varphi_*M)$. A second application of \cite[Prop.~4.1.8]{camargo2026notessolidgeometry} therefore shows that the condition above is equivalent to
\[\varphi_*M\in\mathsf{D}(A)_{\geq 0}\otimes\mathsf{Sp}.\]
This gives the desired description of the induced pre-analytic ring structure.

\textsc{Step IV: $\Psi_\lambda$ is an equivalence on fibers.} We have now established the hypotheses of \cite[Tag~\href{https://kerodon.net/tag/028B}{028B}]{kerodon}, so it remains to show that $\Psi_\lambda$ induces an equivalence on each fiber. For a $\lambda$-light condensed ring $A^\triangleright$, the fiber of $U$ over $A^\triangleright$ is the partially ordered set of pre-analytic ring structures on $A^\triangleright$, ordered by reverse inclusion, and similarly for $U^\mathrm{cn}$. The required equivalence is then a direct consequence of \cite[Thm.~4.4.1]{camargo2026notessolidgeometry}.
\end{proof}

Since presentability is not preserved under passage to opposite categories, it will be convenient to reformulate the definition of $\mathsf{PAnRing}_\lambda^\mathrm{cn}$ in terms of the covariant functor
\[\mathsf{D}(-)_{\geq 0}\colon \mathsf{Cond}_\lambda(\mathsf{Ring})\to \widehat{\mathsf{Cat}},\]
rather than its contravariant counterpart.

\begin{lemma}\label{contractability of fibers}
Let $\varphi\colon A^\triangleright\to B^\triangleright$ be a morphism of $\lambda$-light condensed rings, and consider full subcategories $\mathsf{D}(A)_{\geq0}\subseteq \mathsf{D}(A^\triangleright)_{\geq0}$ and $\mathsf{D}(B)_{\geq0}\subseteq \mathsf{D}(B^\triangleright)_{\geq0}$ closed under small limits and colimits. We retain the notation of Definition~\ref{def connective panring}, but do not assume condition (ii) therein. The following conditions are equivalent:
    \begin{enumerate}[label=(\alph*)]
        \item The restriction-of-scalars $\varphi_*\colon\mathsf{D}(B^\triangleright)_{\geq0} \to\mathsf{D}(A^\triangleright)_{\geq0}$ sends $\mathsf{D}(B)_{\geq0}$ into $\mathsf{D}(A)_{\geq0}$;
        \item There exists a functor $\overline{\varphi}^*\colon \mathsf{D}(A)_{\geq0}\to \mathsf{D}(B)_{\geq0}$ and a natural isomorphism
        \[\alpha_\varphi\colon \bigl(-\otimes_{B^\triangleright}B\bigr)\circ\varphi^*\to \overline{\varphi}^*\circ \bigl(-\otimes_{A^\triangleright}A\bigr)\]
        witnessing the commutativity of the diagram
        % https://q.uiver.app/#q=WzAsNCxbMCwwLCJcXG1hdGhzZntEfShBXlxcdHJpYW5nbGVyaWdodClfe1xcZ2VxMH0iXSxbMCwxLCJcXG1hdGhzZntEfShCXlxcdHJpYW5nbGVyaWdodClfe1xcZ2VxMH0iXSxbMiwwLCJcXG1hdGhzZntEfShBKV97XFxnZXEwfSJdLFsyLDEsIlxcbWF0aHNme0R9KEIpX3tcXGdlcTB9LiJdLFswLDEsIlxcdmFycGhpXioiLDJdLFsxLDMsIi1cXG90aW1lc197Ql5cXHRyaWFuZ2xlcmlnaHR9QiJdLFsyLDMsIlxcb3ZlcmxpbmV7XFx2YXJwaGl9XioiLDAseyJzdHlsZSI6eyJib2R5Ijp7Im5hbWUiOiJkYXNoZWQifX19XSxbMCwyLCItXFxvdGltZXNfe0FeXFx0cmlhbmdsZXJpZ2h0fUEiXV0=
\[\begin{tikzcd}[ampersand replacement=\&]
	{\mathsf{D}(A^\triangleright)_{\geq0}} \&\& {\mathsf{D}(A)_{\geq0}} \\
	{\mathsf{D}(B^\triangleright)_{\geq0}} \&\& {\mathsf{D}(B)_{\geq0}.}
	\arrow["{-\otimes_{A^\triangleright}A}", from=1-1, to=1-3]
	\arrow["{\varphi^*}"', from=1-1, to=2-1]
	\arrow["{\overline{\varphi}^*}", from=1-3, to=2-3]
	\arrow["{-\otimes_{B^\triangleright}B}", from=2-1, to=2-3]
\end{tikzcd}\]
    \end{enumerate}
The anima of pairs $(\overline{\varphi}^*,\alpha_\varphi)$ as in (b) is contractible when these conditions hold, and empty otherwise.
\end{lemma}

\begin{proof}
Assume (a), so that $\varphi_*$ restricts to a functor $\overline{\varphi}_*\colon \mathsf{D}(B)_{\geq0}\to\mathsf{D}(A)_{\geq0}$, and define $\overline{\varphi}^{*}\colonequals
\bigl(-\otimes_{B^\triangleright}B\bigr)\circ\varphi^{*}\circ\iota_A$. For $M\in\mathsf{D}(A)_{\geq0}$ and $N\in\mathsf{D}(B)_{\geq0}$, we have that
\begin{align*}
        \operatorname{Map}_{\mathsf{D}(B)_{\geq0}}(\overline{\varphi}^{*}M,N) &\simeq \operatorname{Map}_{\mathsf{D}(B^{\triangleright})_{\geq0}}(\varphi^{*}\iota_{A}M,\iota_BN) && B\otimes_{B^{\triangleright}}-\dashv \iota_{B} \\
        &\simeq\operatorname{Map}_{\mathsf{D}(A^{\triangleright})_{\geq0}}(\iota_{A}M,\varphi_{*}\iota_B N) && \varphi^{*}\dashv \varphi_{*} \\
        &\simeq\operatorname{Map}_{\mathsf{D}(A^{\triangleright})_{\geq0}}(\iota_{A}M,\iota_A\overline{\varphi}_{*} N) && \varphi_{*}\circ \iota_B\simeq \iota_A\circ \overline{\varphi}_{*} \\
        &\simeq \operatorname{Map}_{\mathsf{D}(A)_{\geq0}}(M,\overline{\varphi}_{*}N) && \iota_A\text{ is fully faithful}.
    \end{align*}
As a result, the functors $\bigl(-\otimes_{B^\triangleright}B\bigr)\circ\varphi^*$ and $\overline{\varphi}^*\circ \bigl(-\otimes_{A^\triangleright}A\bigr)$ are both left adjoint to
\[\varphi_{*}\circ \iota_B\simeq \iota_A\circ \overline{\varphi}_{*},\]
and so they are naturally isomorphic.

Conversely, assume (b). Precomposing $\alpha_\varphi$ with $\iota_A$ and using that $\iota_A$ is fully faithful, we obtain a natural equivalence
\[\overline{\varphi}^*\simeq\overline{\varphi}^*\circ \bigl(-\otimes_{A^\triangleright}A\bigr)\circ\iota_A\simeq\bigl(-\otimes_{B^\triangleright}B\bigr)\circ\varphi^*\circ\iota_A.\]
Since $\iota_A$ preserves small colimits and the other two functors on the right are left adjoints, $\overline{\varphi}^*$ preserves small colimits; hence admits a right adjoint $\overline{\varphi}_*$. Taking right adjoints in $\alpha_\varphi$ yields a natural equivalence
\[\varphi_*\circ\iota_B\simeq\iota_A\circ\overline{\varphi}_*,\]
showing that $\varphi_*$ sends $\mathsf{D}(B)_{\geq0}$ into $\mathsf{D}(A)_{\geq0}$, as required.

Finally, since the completion functor $-\otimes_{A^\triangleright}A\colon \mathsf{D}(A^\triangleright)_{\geq 0}\to \mathsf{D}(A)_{\geq 0}$ is a localization, it induces a fully faithful functor
\[\mathsf{Fun}\mleft(\mathsf{D}(A)_{\geq 0},\mathsf{D}(B)_{\geq 0}\mright)\xhookrightarrow{\hspace{0.8em} -\circ \bigl(-\otimes_{A^\triangleright}A\bigr)\hspace{0.8em} }\mathsf{Fun}\mleft(\mathsf{D}(A^\triangleright)_{\geq 0},\mathsf{D}(B)_{\geq 0}\mright).\]
The induced functor on cores is still fully faithful \cite[Tag~\href{https://kerodon.net/tag/01JP}{01JP}]{kerodon}, and the anima of pairs $(\overline{\varphi}^*,\alpha_\varphi)$ is its fiber over $\bigl(-\otimes_{B^\triangleright}B\bigr)\circ\varphi^*$. It follows that this anima is either empty or contractible \cite[Tags~\href{https://kerodon.net/tag/06SF}{06SF}
and \href{https://kerodon.net/tag/056B}{056B}]{kerodon}; the latter happening precisely when the equivalent conditions are satisfied.
\end{proof}

The following lemma makes precise, with the size bounds needed here, the connective pullback description outlined in \cite[\S2.1.1.9]{Kipp}.

\begin{lemma}\label{PAnRing as a pullback}
Set $\mathsf{W}\colonequals\mathsf{Cond}_\lambda(\mathsf{Ab})$. Then $\mathsf{PAnRing}^\mathrm{cn}_\lambda$ is naturally equivalent to the fiber product of the diagram
\[\mathsf{Cond}_\lambda(\mathsf{Ring})\xrightarrow{\ \mathsf{D}(-)_{\geq 0}\ }\mathsf{CAlg}(\mathsf{Pr}^\mathsf{L}_{\aleph_1})_{\mathsf{W}/}\xleftarrow{\ s\ }\overline{\mathsf{Loc}}_{\aleph_1}^\otimes(\mathsf{W}),\]
where $s$ is the functor sending each symmetric monoidal localization to its source.
\end{lemma}

\begin{proof}
Let $\mathsf{Q}^\mathrm{cn}_\lambda$ denote the fiber product in the statement. Since $\overline{\mathsf{Loc}}^{\otimes}_{\aleph_1}(\mathsf{W})$ is a replete full subcategory of $\mathsf{Fun}(\Delta^{1},\mathsf{CAlg}(\mathsf{Pr}^\mathsf{L}_{\aleph_1})_{\mathsf{W}/})$, the functor
\[s\colon \overline{\mathsf{Loc}}_{\aleph_1}^\otimes(\mathsf{W})\to \mathsf{CAlg}(\mathsf{Pr}^\mathsf{L}_{\aleph_1})_{\mathsf{W}/}\]
is an isofibration by \cite[Tags~\href{https://kerodon.net/tag/01EX}{01EX} and \href{https://kerodon.net/tag/01F3}{01F3}]{kerodon}. Consequently, as at the outset of this section, \cite[Tag~\href{https://kerodon.net/tag/033T}{033T}]{kerodon} identifies the objects of $\mathsf{Q}_{\lambda}^{\mathrm{cn}}$ with pairs $(A^{\triangleright},-\otimes_{A^{\triangleright}}A)$, consisting of a $\lambda$-light condensed ring $A^{\triangleright}$ and a symmetric monoidal localization
\[-\otimes_{A^{\triangleright}}A\colon \mathsf{D}(A^{\triangleright})_{\geq 0}\to\mathsf{D}(A)_{\geq 0}\]
belonging to $\overline{\mathsf{Loc}}^{\otimes}_{\aleph_{1}}(\mathsf{W})$. In particular, it folows from Definition~\ref{def connective panring} -- rephrasing condition (ii) therein via \cite[Lem.~7.5]{reflectionprinciple} -- that $(A^\triangleright,\mathsf{D}(A)_{\geq 0})$ is an object of $\mathsf{PAnRing}^\mathrm{cn}_\lambda$.

The universal property of fiber products yields a functor $\mathsf{Q}^\mathrm{cn}_\lambda\to\mathsf{P}^\mathrm{cn}_\lambda$, induced by the canonical projection $\mathsf{Q}^\mathrm{cn}_\lambda\to\mathsf{Cond}_\lambda(\mathsf{Ring})$ and the composite
\[\overline{\mathsf{Loc}}^{\otimes}_{\kappa}(\mathsf{W})\to\mathsf{Fun}(\Delta^{1},\mathsf{Pr}^{\mathsf{L}}_{\kappa})\simeq\mathsf{Fun}(\Delta^{1},\mathsf{Pr}^{\mathsf{R}}_{\kappa})^{\mathrm{op}}\to\mathsf{Fun}(\Delta^{1},\widehat{\mathsf{Cat}})^{\mathrm{op}},\]
whose essential image is contained in $\mathsf{FF}^{\mathrm{op}}$. As above, the resulting functor takes values in $\mathsf{PAnRing}^\mathrm{cn}_\lambda$ and thus induces a functor
\[\Phi_\lambda\colon\mathsf{Q}^\mathrm{cn}_\lambda\to\mathsf{PAnRing}^\mathrm{cn}_\lambda,\]
which we claim is an equivalence.

To prove that $\Phi_\lambda$ is essentially surjective, let $(A^\triangleright,\mathsf{D}(A)_{\geq 0})$ be an object of $\mathsf{PAnRing}^\mathrm{cn}_\lambda$. As explained in the first step of the proof of Lemma~\ref{preanalytic rings are determined by connective halves}, the completion functor $-\otimes_{A^\triangleright}A$ equips $\mathsf{D}(A)_{\geq 0}$ with a symmetric monoidal structure making
\[-\otimes_{A^\triangleright}A\colon\mathsf{D}(A^\triangleright)_{\geq 0}\to\mathsf{D}(A)_{\geq 0}\]
a morphism in $\mathsf{CAlg}(\mathsf{Pr}^\mathsf{L}_{\aleph_1})$. Equipping $\mathsf{D}(A)_{\geq 0}$ with the induced symmetric monoidal functor from $\mathsf{W}$, we obtain an object $(A^\triangleright,-\otimes_{A^\triangleright}A)$ of $\mathsf{Q}^\mathrm{cn}_\lambda$ satisfying
\[\Phi_\lambda(A^\triangleright,-\otimes_{A^\triangleright}A)\simeq (A^\triangleright,\mathsf{D}(A)_{\geq 0}).\]

To prove that $\Phi_\lambda$ is fully faithful, observe that it is a functor over $\mathsf{Cond}_\lambda(\mathsf{Ring})$. It therefore suffices to show that, for every morphism $\varphi\colon A^\triangleright\to B^\triangleright$ of $\lambda$-light condensed rings, the induced map on fibers over $\varphi$ of
\[\operatorname{Map}_{\mathsf{Q}_{\lambda}^{\mathrm{cn}}}\mleft(\mleft(A^{\triangleright},-\otimes_{A^{\triangleright}}A\mright),\mleft(B^{\triangleright},-\otimes_{B^{\triangleright}}B\mright)\mright)\to\operatorname{Map}_{\mathsf{PAnRing}_\lambda^\mathrm{cn}}\mleft(\mleft(A^\triangleright,\mathsf{D}(A)_{\geq 0}\mright),\mleft(B^\triangleright,\mathsf{D}(B)_{\geq 0}\mright)\mright)\]
is an equivalence. By definition of $\mathsf{PAnRing}_\lambda^\mathrm{cn}$, the fiber on the right is contractible if $\varphi_*$ maps $\mathsf{D}(B)_{\geq 0}$ into $\mathsf{D}(A)_{\geq 0}$, and empty otherwise. Thus, it remains to establish the same description for the fiber on the left.

To establish the same description for the fiber on the left, consider the following cartesian square of mapping animae, induced by the defining fiber product of $\mathsf{Q}^\mathrm{cn}_\lambda$:
% https://q.uiver.app/#q=WzAsNCxbMCwwLCJcXG9wZXJhdG9ybmFtZXtNYXB9X3tcXG1hdGhzZntRfV97XFxsYW1iZGF9XntcXG1hdGhybXtjbn19fVxcbWxlZnQoXFxtbGVmdChBXntcXHRyaWFuZ2xlcmlnaHR9LC1cXG90aW1lc197QV57XFx0cmlhbmdsZXJpZ2h0fX1BXFxtcmlnaHQpLFxcbWxlZnQoQl57XFx0cmlhbmdsZXJpZ2h0fSwtXFxvdGltZXNfe0Jee1xcdHJpYW5nbGVyaWdodH19QlxcbXJpZ2h0KVxcbXJpZ2h0KSJdLFswLDEsIlxcb3BlcmF0b3JuYW1le01hcH1fe1xcbWF0aHNme0NvbmR9X3tcXGxhbWJkYX0oXFxtYXRoc2Z7UmluZyl9fShBXntcXHRyaWFuZ2xlcmlnaHR9LEJee1xcdHJpYW5nbGVyaWdodH0pIl0sWzEsMSwiXFxvcGVyYXRvcm5hbWV7TWFwfV97XFxtYXRoc2Z7Q0FsZ30oXFxtYXRoc2Z7UHJ9XntcXG1hdGhzZntMfX1fe1xcYWxlcGhfezF9fSlfe1xcbWF0aHNme1d9L319XFxtbGVmdChcXG1hdGhzZntEfShBXntcXHRyaWFuZ2xlcmlnaHR9KSxcXG1hdGhzZntEfShCXntcXHRyaWFuZ2xlcmlnaHR9KVxcbXJpZ2h0KSJdLFsxLDAsIlxcb3BlcmF0b3JuYW1le01hcH1fe1xcb3ZlcmxpbmV7XFxtYXRoc2Z7TG9jfX1fe1xcYWxlcGhfezF9fV57XFxvdGltZXN9KFxcbWF0aHNme1d9KX1cXG1sZWZ0KC1cXG90aW1lc197QV57XFx0cmlhbmdsZXJpZ2h0fX1BLC1cXG90aW1lc197Ql57XFx0cmlhbmdsZXJpZ2h0fX1CXFxtcmlnaHQpIl0sWzAsM10sWzMsMiwicyJdLFswLDEsIlUiLDJdLFsxLDJdXQ==
\[\begin{tikzcd}[ampersand replacement=\&]
	{\operatorname{Map}_{\mathsf{Q}_{\lambda}^{\mathrm{cn}}}\mleft(\mleft(A^{\triangleright},-\otimes_{A^{\triangleright}}A\mright),\mleft(B^{\triangleright},-\otimes_{B^{\triangleright}}B\mright)\mright)} \& {\operatorname{Map}_{\overline{\mathsf{Loc}}_{\aleph_{1}}^{\otimes}(\mathsf{W})}\mleft(-\otimes_{A^{\triangleright}}A,-\otimes_{B^{\triangleright}}B\mright)} \\
	{\operatorname{Map}_{\mathsf{Cond}_{\lambda}(\mathsf{Ring)}}(A^{\triangleright},B^{\triangleright})} \& {\operatorname{Map}_{\mathsf{CAlg}(\mathsf{Pr}^{\mathsf{L}}_{\aleph_{1}})_{\mathsf{W}/}}\mleft(\mathsf{D}(A^{\triangleright})_{\geq 0},\mathsf{D}(B^{\triangleright})_{\geq 0}\mright).}
	\arrow[from=1-1, to=1-2]
	\arrow["U"', from=1-1, to=2-1]
	\arrow["s", from=1-2, to=2-2]
	\arrow[from=2-1, to=2-2]
\end{tikzcd}\]
The pasting law gives an equivalence between the fiber of $U$ over $\varphi$ and the fiber of $s$ over $\varphi^*$. By Lemma~\ref{contractability of fibers}, the latter is contractible if $\varphi_*$ maps $\mathsf{D}(B)_{\geq 0}$ into $\mathsf{D}(A)_{\geq 0}$, and empty otherwise. This establishes the desired description and completes the proof.
\end{proof}

\begin{proof}[Proof of Theorem~\ref{PAnRing is presentable}]
Lemmas~\ref{preanalytic rings are determined by connective halves} and \ref{PAnRing as a pullback} identify $\mathsf{PAnRing}_\lambda$ with the pullback of the diagram
\[\mathsf{Cond}_\lambda(\mathsf{Ring})\xrightarrow{\ \mathsf{D}(-)_{\geq 0}\ }\mathsf{CAlg}(\mathsf{Pr}^\mathsf{L}_{\aleph_1})_{\mathsf{W}/}\xleftarrow{\ s\ }\overline{\mathsf{Loc}}^\otimes_{\aleph_1}(\mathsf{W}),\]
where $\mathsf{W}=\mathsf{Cond}_\lambda(\mathsf{Ab})$. By Proposition~\ref{D is a map in Pr_k}, the morphism on the left lies in $\mathsf{Pr}^\mathsf{L}_{\aleph_1}$. On the other hand, Proposition~\ref{CondAb is mu-compact} allows us to apply Corollary~\ref{main theorem on monoidal localizations2}, which shows that the morphism on the right lies in $\mathsf{Acc}_\mu$ and preserves $\aleph_1$-filtered colimits.

The argument of \cite[Tag~\href{https://kerodon.net/tag/06LM}{06LM}]{kerodon} therefore shows that $\mathsf{PAnRing}_\lambda$ is $\mu$-accessible. Finally, \cite[Lem.~4.1.10 and Prop.~4.1.11]{camargo2026notessolidgeometry} proves that $\mathsf{PAnRing}_{\aleph_1}$ admits small colimits. The same arguments apply verbatim for every uncountable cardinal $\lambda$, and hence $\mathsf{PAnRing}_\lambda$ admits small colimits in general.
\end{proof}

A convenient feature of presentable categories is that they are automatically complete \cite[Tag~\href{https://kerodon.net/tag/06PU}{06PU}]{kerodon}. We therefore obtain the following corollary, for which we are not aware of a direct proof.

\begin{corollary}\label{PAnRing is complete}
   The category $\mathsf{PAnRing}_\lambda$ is complete. 
\end{corollary}

In Clausen--Scholze's theory of analytic geometry, one typically works with \emph{analytic rings}, namely pre-analytic rings $A$ satisfying $A^\triangleright\in\mathsf{D}(A)$. In the light setting, where $\lambda=\aleph_1$, pre-analytic rings enjoy an additional important property: for every prime $p$, the Frobenius morphism
\[\phi_p\colon A^\triangleright\to A^\triangleright/p\]
underlies a morphism of pre-analytic rings $A\to A/p$, where $A/p$ carries the pre-analytic ring structure induced by $A$ \cite[Thm.~4.3.2]{camargo2026notessolidgeometry}. As it is unclear whether this property extends to arbitrary $\lambda$, we henceforth restrict to the light setting and suppress $\lambda$ from the notation.

\begin{definition}
Let $A$ be a light pre-analytic ring. We say that $A$ is an \emph{analytic ring} if the object $A^\triangleright$ of $\mathsf{D}(A^\triangleright)$ belongs to the full subcategory $\mathsf{D}(A)$. We denote by $\mathsf{AnRing}\subseteq\mathsf{PAnRing}$ the full subcategory spanned by the analytic rings.
\end{definition}

We are now in a position to state the main application of the machinery developed in Sections~\ref{sec2} and \ref{sec3}. As noted in the introduction, the following theorem was announced in \cite[Lecture~15]{ClausenScholzeAnalyticStacks} without proof and raised as an open point in \cite[\S1.2, Footnote~2]{camargo2026notessolidgeometry}.

\begin{theorem}\label{AnRing is presentable}
    The category $\mathsf{AnRing}$ is $\mu$-presentable, for $\mu=(2^{\aleph_0})^+$.
\end{theorem}

\begin{proof}
According to \cite[Thm.~4.1.13 and Cor.~4.2.7]{camargo2026notessolidgeometry}, the natural inclusion
\[\mathsf{AnRing}\hookrightarrow\mathsf{PAnRing}\]
admits a left adjoint. It follows from Corollary~\ref{PAnRing is complete} and \cite[Tag~\href{https://kerodon.net/tag/04JX}{04JX}]{kerodon} that $\mathsf{AnRing}$ is complete and that the inclusion preserves and reflects small limits. By \cite[Cor.~4.1.12]{camargo2026notessolidgeometry}, we also have that the category $\mathsf{AnRing}$ is stable under sifted, hence filtered, colimits in $\mathsf{PAnRing}$. The desired result then follows from Theorem~\ref{PAnRing is presentable} and Proposition~\ref{enhanced reflection theorem}.
\end{proof}

For convenience, we record the following consequence of the preceding proof.

\begin{corollary}\label{AnRing is complete}
   The category $\mathsf{AnRing}$ is complete. Moreover, the inclusion
   \[\mathsf{AnRing}\hookrightarrow\mathsf{PAnRing}\] 
   preserves and reflects small limits. 
\end{corollary}